\section*{अध्याय \ref{ch:g3:separating-mixtures} --- मिश्रणों को अलग करना}

\begin{solution}{exo:g3:separating-mixtures:1}
\omterm{def:g3:separating-mixtures:sieve}{चालना}: चावल के दाने मटर से छोटे हैं और सही छेदों वाली छलनी से नीचे
गिर जाते हैं। (हाथ से छाँटना भी काम करता है, पर धीरे-धीरे।)
\end{solution}

\begin{solution}{exo:g3:separating-mixtures:2}
मिट्टी धीरे-धीरे नीचे जाकर तली में परत बना लेती है; ऊपर का पानी
ज़्यादा स्वच्छ हो जाता है। इसे \omterm{def:g3:separating-mixtures:decant}{अवसादन} कहते हैं।
\end{solution}

\begin{solution}{exo:g3:separating-mixtures:3}
\omterm{def:g3:separating-mixtures:filter}{निस्यंद} वह द्रव है जो छन्ने से छनकर निकलता है। कॉफ़ी बनाते समय
\omterm{def:g3:separating-mixtures:filter}{निस्यंद} जग में आई कॉफ़ी है; पिसी कॉफ़ी कागज़ में रह
जाती है।
\end{solution}

\begin{solution}{exo:g3:separating-mixtures:4}
हाँ। नमक घुला हुआ है, और छन्ना घुली हुई चीज़ को नहीं रोकता: वह पानी के
साथ निकल जाता है।
\end{solution}

\begin{solution}{exo:g3:separating-mixtures:5}
पानी का \omterm{def:g3:separating-mixtures:evaporate}{वाष्पन} करके: धूप में, या हल्का गरम करने पर, पानी हवा में चला
जाता है और चीनी प्याली में रह जाती है।
\end{solution}

\begin{solution}{exo:g3:separating-mixtures:6}
बजरी छलनी में रह जाती है। रेत के कण छेदों से छोटे हैं, इसलिए वे
नीचे गिर जाते हैं।
\end{solution}

\begin{solution}{exo:g3:separating-mixtures:7}
(क) \omterm{def:g3:separating-mixtures:sieve}{चालना}। (ख) \omterm{def:g3:separating-mixtures:filter}{छानना} (या \omterm{def:g3:separating-mixtures:decant}{अवसादन} और \omterm{def:g3:separating-mixtures:decant}{निस्तारण})। (ग) \omterm{def:g3:separating-mixtures:evaporate}{वाष्पन}।
\end{solution}

\begin{solution}{exo:g3:separating-mixtures:8}
चार चरण: जालियाँ, \omterm{def:g3:separating-mixtures:decant}{अवसादन} टंकियाँ, रेत के छन्ने, कीटाणुओं का नाश।
\omterm{def:g3:separating-mixtures:decant}{अवसादन} टंकियाँ कीचड़ हटाती हैं।
\end{solution}

\begin{solution}{exo:g3:separating-mixtures:9}
छानो; पेपर से रेत इकट्ठा करो; नमक पाने के लिए \omterm{def:g3:separating-mixtures:filter}{निस्यंद} का
\omterm{def:g3:separating-mixtures:evaporate}{वाष्पन} करो।
\end{solution}

\begin{solution}{exo:g3:separating-mixtures:10}
नहीं। समुद्र के पानी का नमक घुला हुआ है, और छन्ना घुली हुई चीज़ को नहीं
रोकता: \omterm{def:g3:separating-mixtures:filter}{निस्यंद} अब भी नमकीन पानी ही है।
\end{solution}

\begin{solution}{exo:g3:separating-mixtures:11}
पत्तियों के नीचे बैठ जाने पर चाय को धीरे से अलग उँडेलना (\omterm{def:g3:separating-mixtures:decant}{निस्तारण}), या
उसे छलनी या छन्ने से उँडेलना (\omterm{def:g3:separating-mixtures:sieve}{चालना} या \omterm{def:g3:separating-mixtures:filter}{छानना})।
\end{solution}
